
\documentclass[12pt,a4paper]{article}

% Paquetes
\usepackage[utf8]{inputenc}
\usepackage[spanish]{babel}
\usepackage{geometry}
\geometry{margin=2.5cm}
\usepackage{graphicx}
\usepackage{booktabs}
\usepackage{longtable}
\usepackage{array}
\usepackage{hyperref}
\usepackage{xcolor}
\usepackage{fancyhdr}
\usepackage{titlesec}
\usepackage{amsmath}
\usepackage{pgfplots}
\pgfplotsset{compat=1.18}
\usepackage{float}
\usepackage{caption}
\usepackage{subcaption}

% Colores personalizados
\definecolor{ustblue}{RGB}{0,51,102}
\definecolor{ustgray}{RGB}{128,128,128}
\definecolor{sigreen}{RGB}{39,174,96}
\definecolor{sired}{RGB}{231,76,60}
\definecolor{siyellow}{RGB}{241,196,15}

% Configuración de hipervínculos
\hypersetup{
    colorlinks=true,
    linkcolor=ustblue,
    urlcolor=ustblue,
    citecolor=ustblue
}

% Encabezados y pies de página
\pagestyle{fancy}
\fancyhf{}
\fancyhead[L]{\textcolor{ustgray}{Análisis de Bienestar Psicológico}}
\fancyhead[R]{\textcolor{ustgray}{\today}}
\fancyfoot[C]{\thepage}
\renewcommand{\headrulewidth}{0.4pt}
\renewcommand{\footrulewidth}{0.4pt}

% Títulos
\titleformat{\section}{\Large\bfseries\color{ustblue}}{\thesection}{1em}{}
\titleformat{\subsection}{\large\bfseries\color{ustblue!80}}{\thesubsection}{1em}{}
\titleformat{\subsubsection}{\normalsize\bfseries\color{ustblue!60}}{\thesubsubsection}{1em}{}

\begin{document}

% Portada
\begin{titlepage}
\centering
\vspace*{2cm}

{\Huge\bfseries\textcolor{ustblue}{Análisis de Bienestar Psicológico Estudiantil}\par}

\vspace{1cm}

{\Large\textcolor{ustgray}{Universidad Santo Tomás}\par}

\vspace{0.5cm}

{\large\textcolor{ustgray}{Universidad Santo Tomás}\par}

\vspace{2cm}

{\large\textcolor{ustgray}{Fecha: 25 de August de 2026}\par}

\vspace{1cm}

\begin{tikzpicture}
\draw[fill=ustblue!10, draw=ustblue, rounded corners=5pt] (0,0) rectangle (12,4);
\node[align=center] at (6,3) {\textbf{\large Resumen Ejecutivo}};
\node[align=center] at (6,2) {Muestra: 281 estudiantes};
\node[align=center] at (6,1.2) {Bienestar Global: 4.52 $\pm$ 0.65};
\end{tikzpicture}

\vfill

{\small\textcolor{ustgray}{Generado automáticamente por el Sistema de Análisis de Bienestar Estudiantil}\par}
\end{titlepage}

% Índice
\tableofcontents
\newpage

% ============================================================
% SECCIÓN 1: INTRODUCCIÓN
% ============================================================
\section{Introducción}

\subsection{Contexto}
El bienestar psicológico de los estudiantes universitarios es un factor determinante 
en su rendimiento académico, retención y desarrollo integral. Este informe presenta 
un análisis completo de 281 respuestas obtenidas a través 
de la Escala de Bienestar Psicológico de Ryff.

\subsection{Objetivos}
\begin{itemize}
    \item Evaluar el nivel de bienestar psicológico de los estudiantes
    \item Identificar dimensiones con mayores áreas de oportunidad
    \item Detectar patrones y perfiles de estudiantes
    \item Proponer intervenciones basadas en evidencia
\end{itemize}

\subsection{Metodología}
Se utilizó la Escala de Bienestar Psicológico de Ryff (29 ítems), adaptada al español 
por Díaz et al. (2006). Las respuestas se codificaron en una escala Likert de 1 a 6, 
donde valores más altos indican mayor bienestar.

% ============================================================
% SECCIÓN 2: RESULTADOS
% ============================================================
\section{Resultados}

\subsection{Estadísticas Descriptivas}

\begin{table}[H]
\centering
\caption{Estadísticas descriptivas por dimensión del bienestar}
\label{tab:descriptivas}
\begin{tabular}{@{}lcccccc@{}}
\toprule
\textbf{Dimensión} & \textbf{N} & \textbf{Media} & \textbf{DE} & \textbf{Mín} & \textbf{Máx} & \textbf{Interpretación} \\
\midrule
Autoaceptación & 268 & 4.72 & 0.93 & 1 & 6 & \textcolor{siyellow}{Medio} \\
Relaciones Positivas & 268 & 4.28 & 0.89 & 2 & 6 & \textcolor{siyellow}{Medio} \\
Autonomía & 268 & 4.34 & 0.84 & 2 & 6 & \textcolor{siyellow}{Medio} \\
Dominio del Entorno & 268 & 4.35 & 0.78 & 2 & 6 & \textcolor{siyellow}{Medio} \\
Propósito de Vida & 268 & 4.66 & 0.81 & 2 & 6 & \textcolor{siyellow}{Medio} \\
Crecimiento Personal & 268 & 4.79 & 0.79 & 2 & 6 & \textcolor{siyellow}{Medio} \\

\bottomrule
\end{tabular}
\end{table}

\subsection{Índice Global de Bienestar}

El bienestar global promedio es de \textbf{4.52} sobre 6.00 
(Desviación Estándar = 0.65).

\begin{figure}[H]
\centering
\begin{tikzpicture}
\begin{axis}[
    ybar,
    bar width=15pt,
    width=0.8\textwidth,
    height=8cm,
    ylabel={Puntuación Promedio},
    symbolic x coords={Aut, Rel, Aut, Dom, Pro, Cre},
    xtick=data,
    ymin=0, ymax=6,
    nodes near coords,
    nodes near coords align={vertical},
]
\addplot[fill=ustblue] coordinates {
(Aut, 4.72)
(Rel, 4.28)
(Aut, 4.34)
(Dom, 4.35)
(Pro, 4.66)
(Cre, 4.79)

};
\end{axis}
\end{tikzpicture}
\caption{Puntuaciones promedio por dimensión del bienestar}
\label{fig:dimensiones}
\end{figure}


\subsection{Análisis de Correlaciones}

\begin{table}[H]
\centering
\caption{Correlaciones con el bienestar global}
\label{tab:correlaciones}
\begin{tabular}{@{}lcc@{}}
\toprule
\textbf{Dimensión} & \textbf{Correlación} & \textbf{p-valor} \\
\midrule
Autoaceptación & 0.8166 & 0.0000 *** \\
Relaciones Positivas & 0.6439 & 0.0000 *** \\
Autonomía & 0.7351 & 0.0000 *** \\
Dominio del Entorno & 0.8350 & 0.0000 *** \\
Propósito de Vida & 0.8389 & 0.0000 *** \\
Crecimiento Personal & 0.6874 & 0.0000 *** \\

\bottomrule
\end{tabular}
\end{table}


\subsection{Análisis de Regresión}

\subsubsection{Regresión Lineal Múltiple}
El modelo de regresión lineal múltiple explicó el \textbf{100.0\%} 
de la varianza en el bienestar global (R² = 1.0000).

\begin{table}[H]
\centering
\caption{Coeficientes de regresión estandarizados}
\label{tab:regresion}
\begin{tabular}{@{}lccc@{}}
\toprule
\textbf{Variable} & \textbf{β} & \textbf{p-valor} & \textbf{Significativo} \\
\midrule
Autoaceptación & 0.1546 & 0.0000 & Sí \\
Relaciones Positivas & 0.1487 & 0.0000 & Sí \\
Autonomía & 0.1394 & 0.0000 & Sí \\
Propósito de Vida & 0.1345 & 0.0000 & Sí \\
Crecimiento Personal & 0.1319 & 0.0000 & Sí \\
Dominio del Entorno & 0.1299 & 0.0000 & Sí \\

\bottomrule
\end{tabular}
\end{table}


\subsection{Análisis de Clustering}

Se identificaron \textbf{2} perfiles diferenciados de estudiantes:

\begin{table}[H]
\centering
\caption{Perfiles de estudiantes identificados}
\label{tab:clusters}
\begin{tabular}{@{}lccc@{}}
\toprule
\textbf{Perfil} & \textbf{Cantidad} & \textbf{Porcentaje} & \textbf{Descripción} \\
\midrule
Perfil 0 & 137 & 48.8\% & Autoaceptación - Crecimiento Personal \\
Perfil 1 & 131 & 46.6\% & Crecimiento Personal - Autoaceptación \\

\bottomrule
\end{tabular}
\end{table}


\subsubsection{Regresión Logística}

El modelo logístico clasificó a los estudiantes con \textbf{100.0\%} 
de precisión (AUC = 1.0000).

\textbf{Odds Ratios:}

\begin{itemize}
    \item Autoaceptación: OR = 0.1821 (reduce la probabilidad)
    \item Relaciones Positivas: OR = 0.4373 (reduce la probabilidad)
    \item Autonomía: OR = 0.2553 (reduce la probabilidad)
    \item Dominio del Entorno: OR = 0.4991 (reduce la probabilidad)
    \item Propósito de Vida: OR = 0.2379 (reduce la probabilidad)
    \item Crecimiento Personal: OR = 0.3382 (reduce la probabilidad)

\end{itemize}


\subsubsection{Selección de Variables (Stepwise)}

El análisis stepwise seleccionó \textbf{6} 
variables como predictoras significativas:

\begin{itemize}
    \item score_PV
    \item score_PR
    \item score_CP
    \item score_AU
    \item score_SA
    \item score_DE

\end{itemize}


% ============================================================
% SECCIÓN 3: CONCLUSIONES
% ============================================================
\section{Conclusiones y Recomendaciones}

\subsection{Principales Hallazgos}

\begin{itemize}
    \item El bienestar global es \textcolor{siyellow}{\textbf{MEDIO}}, existiendo 
    oportunidades de mejora en varias dimensiones.
    \item Se recomienda implementar programas focalizados en las dimensiones más débiles.
\end{itemize}

\subsection{Recomendaciones Específicas}

Basado en el análisis, la dimensión con mayor área de oportunidad es 
\textbf{Relaciones Positivas} (Media = 4.28).

\begin{enumerate}
    \item \textbf{Diagnóstico profundo}: Realizar grupos focales para entender las 
    causas raíz de esta dimensión.
    \item \textbf{Intervención diseñada}: Crear programas específicos para mejorar 
    esta dimensión.
    \item \textbf{Seguimiento}: Implementar mediciones periódicas para evaluar el impacto.
    \item \textbf{Capacitación}: Entrenar al personal docente en estrategias de apoyo.
\end{enumerate}

% ============================================================
% ANEXOS
% ============================================================
\newpage
\appendix
\section{Anexos}

\subsection{Metodología Detallada}
La Escala de Bienestar Psicológico de Ryff (1989) evalúa seis dimensiones:
\begin{enumerate}
    \item \textbf{Autoaceptación}: Actitud positiva hacia uno mismo
    \item \textbf{Relaciones Positivas}: Conexiones cálidas con otros
    \item \textbf{Autonomía}: Independencia y auto-regulación
    \item \textbf{Dominio del Entorno}: Gestión efectiva del entorno
    \item \textbf{Propósito de Vida}: Sentido y dirección en la vida
    \item \textbf{Crecimiento Personal}: Desarrollo continuo
\end{enumerate}

\subsection{Ítems Invertidos}
Los siguientes ítems fueron recodificados: 2, 4, 5, 8, 9, 13, 19, 22, 23, 26.

\vfill

\begin{center}
\textcolor{ustgray}{\small Documento generado automáticamente el 25/08/2026 19:03}
\end{center}

\end{document}
